% TOP_PROBLEM: {"id": 30006745, "title": "The Generalized Sato-Tate Conjecture for Higher-Dimensional Abelian Varieties", "category_id": 1, "category": {"id": 1, "name": "number_theory", "display_name": "Number Theory", "description": "Properties of integers, prime numbers, Diophantine equations.", "slug": "number-theory", "order_index": 1, "created_at": "2026-07-31T15:26:25.670Z"}, "status": "open"}
% FIRST_POSTED: null
% !TeX program = lualatex
% Compile twice: lualatex 30006745.tex
% All references and prose are embedded; no BibTeX database or external inputs.
\documentclass[11pt,a4paper]{article}
\usepackage[margin=24mm,headheight=14pt]{geometry}
\usepackage{fontspec}
\setmainfont{Latin Modern Roman}
\setsansfont{Latin Modern Sans}
\setmonofont{Latin Modern Mono}
\usepackage{amsmath,amssymb,mathtools}
\usepackage{unicode-math}
\setmathfont{Latin Modern Math}
\usepackage{microtype}
\usepackage{enumitem,xurl,needspace}
\usepackage[dvipsnames]{xcolor}
\usepackage{hyperref}
\hypersetup{unicode=true,colorlinks=true,linkcolor=MidnightBlue,urlcolor=MidnightBlue,
 pdftitle={Research attempt: Problem 30006745},pdfauthor={Research notebook},bookmarksnumbered=true}
\usepackage{bookmark}
\usepackage{tocloft}
\setlength{\cftsecnumwidth}{3em}
\cftsetpnumwidth{2.5em}
\cftsetrmarg{4em}
\usepackage{titlesec}
\titleformat{\section}{\Large\bfseries\raggedright}{\thesection}{0.7em}{}
\usepackage{fancyhdr}
\pagestyle{fancy}\fancyhf{}
\fancyhead[L]{\small\sffamily Open Problems Math | Research notebook}
\fancyhead[R]{\small Problem 30006745 | 27 September 2026}
\fancyfoot[C]{\thepage}
\setlength{\parindent}{0pt}
\setlength{\parskip}{3pt plus 1pt}
\setlength{\emergencystretch}{3em}
\setcounter{tocdepth}{1}
\setcounter{secnumdepth}{1}
\setlist[itemize]{leftmargin=3em,itemsep=5pt,topsep=4pt,parsep=0pt}
\allowdisplaybreaks
\newcommand{\N}{\mathbb{N}}
\newcommand{\Z}{\mathbb{Z}}
\newcommand{\Q}{\mathbb{Q}}
\newcommand{\R}{\mathbb{R}}
\newcommand{\C}{\mathbb{C}}
\newcommand{\F}{\mathbb{F}}
\newcommand{\A}{\mathbb{A}}
\newcommand{\PP}{\mathbb P}
\newcommand{\E}{\mathbb{E}}
\newcommand{\eps}{\varepsilon}
\newcommand{\dd}{\,\mathrm d}
\newcommand{\sref}[2]{\hyperlink{source:#1:#2}{[S#2]}}
\newcommand{\eref}[2]{\hyperlink{extra:#1:#2}{[E#2]}}
% Turn the next switch off for a shorter reading copy. The quotations preserve
% every exactTarget field, including qualifications omitted from a short summary.
\newif\ifshowcatalogue
\showcataloguetrue
\newcommand{\cataloguescope}[1]{\ifshowcatalogue
 \paragraph{Exact scope in the supplied catalogue.}
 \begin{quote}\small #1\end{quote}\fi}
\usepackage{etoolbox}
\AtBeginEnvironment{itemize}{\small}
\numberwithin{equation}{section}
% Standalone export. Original research content preserved.
\begin{document}
\setcounter{section}{0}
% ================================================================
% problemId: 30006745
% Canonical problem ID: 30006745
% exactTarget SHA-256: 2e1b8ce0b50b58b796ea6a6143f9a8ab54e30fb7eaa256e237e57b56b67e5410
\clearpage
\section*{\texorpdfstring{The Generalized Sato-Tate Conjecture for Higher-Dimensional Abelian Varieties}{The Generalized Sato-Tate Conjecture for Higher-Dimensional Abelian Varieties}}\label{30006745}
\subsection{Definitions and mathematical statement}
For an abelian variety $A/K$ of dimension $g$ and a good finite place $v$, Frobenius on its first cohomology has eigenvalues of absolute value $(Nv)^{1/2}$. Divide them by this factor and use the resulting conjugacy class $c_v$ in the source's compact Sato--Tate group $\operatorname{ST}(A)$. The conjecture is
\[
 \frac1{\pi_K(x)}\sum_{Nv\le x,\ v\text{ good}}F(c_v)
 \longrightarrow\int_{\operatorname{ST}(A)}F(g)\,dg
\]
for every continuous class function $F$, with normalized Haar measure and $\pi_K(x)$ the number of prime ideals of norm at most $x$. The full group, including its components, and its Frobenius-class construction use the source convention.
\par\smallskip\textit{Formulation and concept references:} \sref{30006745}{1}.
\cataloguescope{Prove that the normalized Frobenius conjugacy classes of an abelian variety A of dimension g over a number field K are equidistributed in the Sato-Tate group ST(A).}
\subsection{Short English statement}
Are normalized Frobenius symmetries of every abelian variety distributed according to Haar measure on its Sato--Tate group?
% BEGIN RESEARCH ADDITION 135
\subsection{Research attempt}

\paragraph{Editorial note on the formulation reference.}
The DOI in [S1] resolves to Gan--Takeda, \emph{The local Langlands conjecture
for $\mathrm{GSp}(4)$}, not the displayed authors or a definition of the
Sato--Tate group for arbitrary abelian varieties. The original reference and
statement are retained. In this attempt fix $\ell$ and an embedding
$\overline\Q_\ell\hookrightarrow\C$ and use Sutherland's compact-group and
normalized-Frobenius construction, Sections 3.2--3.3, Conjecture 3.6.
This is a declared convention where [S1] supplies none; no independence of
all choices is assumed. \eref{30006745}{1}\eref{30006745}{2}

\paragraph{Current frontier.}
Sato--Tate for non-CM Hilbert modular forms is an established specialized
result, not the arbitrary higher-dimensional statement. \eref{30006745}{3}
Classification of Sato--Tate groups of abelian threefolds, examples and
moment computations likewise do not establish equidistribution for every
variety with one of those groups. \eref{30006745}{4} The analytic route requires
control of the nontrivial representation $L$-functions, not just knowledge
of the algebraic monodromy group. \eref{30006745}{2}

\paragraph{Attack route.}
Combine a tensor description of the compact group with prime averages.
Possible bridges are automorphy and boundary nonvanishing for all irreducible
$L$-functions, direct cancellation of projected tensor traces, or finite
moment reconstruction. We prove the exact sufficiency of the second bridge
and test the third. The deductions do not claim a new arithmetic
cancellation theorem.

\paragraph{Attempt: projected tensor moments suffice.}
Let $G\subset U(V)$ be compact with faithful $V$, and put
\[
 T_{a,b}=V^{\otimes a}\otimes\overline V^{\otimes b},\qquad
 \Pi_{a,b}=\int_G T_{a,b}(g)\,dg.
\]
This finite-dimensional integral is the orthogonal projection onto invariant
vectors. For every $G$-equivariant orthogonal projection $P$ on $T_{a,b}$,
\[
 F_P(g)=\operatorname{tr}(P T_{a,b}(g))
\]
is a continuous class function with Haar mean $\operatorname{tr}(P\Pi_{a,b})$.
We claim that the prime-average condition, for every fixed $a,b,P$,
\begin{equation}\label{30006745:moments}
 \frac1{\pi_K(x)}\sum_{Nv\le x}F_P(c_v)
       \longrightarrow\operatorname{tr}(P\Pi_{a,b})
\end{equation}
implies the target equidistribution for this group. Finitely many omitted
bad primes have no effect on the normalized averages.

Matrix coefficients of $V$ and $\overline V$ generate a self-adjoint algebra
separating elements of $G$, hence are uniformly dense in $C(G)$ by
Stone--Weierstrass. Products of coefficients are coefficients of the tensors
$T_{a,b}$. If an irreducible representation $W$ occurred in none of these,
Schur orthogonality would make its nonzero coefficients orthogonal to a
dense subspace of $L^2(G)$, a contradiction. Thus some tensor contains $W$.
Projection onto one invariant copy of $W$ commutes with $G$ and its $F_P$
is $\chi_W$. Condition~\eqref{30006745:moments} gives average zero for each
nontrivial irreducible character and one for the trivial character.

To complete the limiting step, conjugation-average a uniformly approximating
coefficient polynomial. Averaging contracts the sup norm; Schur's lemma
makes the conjugation average of an irreducible matrix coefficient a scalar
multiple of its character. Finite character combinations are consequently
dense in the continuous class functions. Approximate a test within
$\epsilon$, use the finitely many character limits, then send
$\epsilon\downarrow0$. No infinite series and prime limit have been
interchanged. This proves the claim.

\paragraph{Why the projections carry essential information.}
Taking only $P=I$ measures powers of $\operatorname{tr}V$ and its conjugate.
Trace need not separate conjugacy classes. For
$G=SU(2)\times SU(2)$ and $V=V_1\oplus V_2$, swapping two different component
traces preserves the total trace but changes the $G$-conjugacy class.
Even all characteristic-polynomial coefficients distinguish only ambient
unitary conjugacy, which can be coarser than conjugacy in a subgroup.
The individual equivariant projections in \eqref{30006745:moments} avoid that
loss. They are not claimed to be easy arithmetic observables.

The group also matters for numerical moments. For diagonal $SU(2)$ on
$V_1\oplus V_1$, the trace is $2t$ and its second Haar moment is $4$.
For the product group on $V_1\oplus V_2$ it is $t_1+t_2$ and the second
moment is $2$. Indeed $\int t=0$ and $\int t^2=1$, by the absence of
invariants in $V_1$ and the one-dimensional invariant subspace of
$V_1^{\otimes2}$. Replacing the trace law of $E^2$ by two independent
elliptic trace laws would thus give the wrong moment even if each marginal
were known perfectly.

\paragraph{Stress test: every fixed truncation can be fooled.}
On $U(1)$, for $D\ge0$ and $0<\epsilon<1$, the probability measure
\[
 \frac{1+\epsilon\cos((D+1)\theta)}{2\pi}\,d\theta
\]
is non-Haar and positive, but agrees with Haar on
$z^a\overline z^b$ for all $a+b\le D$. Orthogonality of integer
exponentials proves the identity exactly. Thus a finite moment test is
insufficient even without numerical error. This constructs probability
measures, not a counterexample Frobenius sequence.

Passing to a field where the group becomes connected does not by itself
prove all averages on the original disconnected group. The component-field
construction is described in Sutherland, Theorem 3.12. \eref{30006745}{2}
The criterion above retains representations of the whole group, including
its finite quotient. Finally, tensor constructions can have repeated
Hodge--Tate weights. Regularity hypotheses in potential automorphy theorems
must be checked, not assumed for every tensor factor. No required
continuation or nonvanishing theorem for all those $L$-functions is proved
here.

\paragraph{Outcome.}
\textbf{Outcome: Conditional reduction.}
All projected tensor prime averages in \eqref{30006745:moments} imply the full
class equidistribution target, with the declared group convention. The
proof retains component information and handles approximation uniformly.
Finite trace moments, independent-product heuristics and connected-component
arguments do not supply the hypothesis. The missing assertion remains
arithmetic cancellation for all the required projected traces.

% END RESEARCH ADDITION 135
\subsection{Sources}
\begin{itemize}\raggedright
\item[\textnormal{[S1]}]\hypertarget{source:30006745:1}{} T. Barnet-Lamb, D. Geraghty, M. Harris, R. Taylor, Ann. of Math. 173(3):1841-1869, 2011.
\url{https://doi.org/10.4007/annals.2011.173.3.12}.
{\small\textit{Catalogue formulation source.}}
% BEGIN RESEARCH REFERENCES 135
\item[\textnormal{[E1]}]\hypertarget{extra:30006745:1}{} W. T. Gan and S. Takeda,
\emph{The local Langlands conjecture for $\mathrm{GSp}(4)$}, Annals of
Mathematics 173 (2011), 1841--1882; metadata for the DOI in [S1].
\url{https://annals.math.princeton.edu/2011/173-3/p12}.
\item[\textnormal{[E2]}]\hypertarget{extra:30006745:2}{} A. V. Sutherland,
\emph{Sato--Tate distributions}, Contemporary Mathematics 740 (2019),
197--248, Sections 2, 3.2--3.3, Conjecture 3.6, Theorem 3.12;
arXiv:1604.01256. \url{https://math.mit.edu/~drew/conm14904.pdf}.
\item[\textnormal{[E3]}]\hypertarget{extra:30006745:3}{} T. Barnet-Lamb, T. Gee
and D. Geraghty, \emph{The Sato--Tate conjecture for Hilbert modular forms},
JAMS 24 (2011), 411--469; arXiv:0912.1054, main theorem.
\url{https://arxiv.org/abs/0912.1054}.
\item[\textnormal{[E4]}]\hypertarget{extra:30006745:4}{} F. Fit\'e, K. S. Kedlaya
and A. V. Sutherland, \emph{Sato--Tate groups of abelian threefolds},
arXiv:2106.13759, abstract and classification statement.
\url{https://arxiv.org/abs/2106.13759}.

% END RESEARCH REFERENCES 135
\end{itemize}

\end{document}
